\chapter{From path integrals to committed quotes}\label{ch:quotes}
A derivative tells us the effect of an infinitesimal displacement. A finite
action travels through changing derivatives. Integrating them recovers the
endpoint difference. This is the mathematical connection between force and
exact EBU, and it remains valid whether or not a controller uses the force
to choose the action.

\section{One action, one chain rule}
Let $\gamma(s)=z+s\delta$ for $0\leq s\leq1$, with $V$ continuously
differentiable along this path. The chain rule gives
\begin{equation}
 \frac{d}{ds}V(\gamma(s))=\nabla V(\gamma(s))^T\delta.
\end{equation}
Integrate and change the sign:
\begin{equation}\label{eq:path}
 V(z)-V(z+\delta)
 =-\int_0^1\nabla V(z+s\delta)^T\delta\,ds.
\end{equation}
For $\delta=qS_e$ with $q\geq0$, substitution $r=sq$ yields
$\int_0^q f_e(z+rS_e)\,dr$. The changing force, not just its initial
value, is integrated. A declared process burden is then subtracted separately.

For the quadratic, $\mu(z+s\delta)=\mu(z)+sH\delta$. Integrating one
and $s$ over $[0,1]$ immediately gives the finite expansion of Chapter 37.
The average endpoint gradient is also exact:
\begin{equation}
 V(z)-V(z+\delta)
 =-\tfrac12[\mu(z)+\mu(z+\delta)]^T\delta.
\end{equation}
This is not a general exact quadrature claim for arbitrary potentials. It
works here because the gradient along the segment is affine.

\section{Path arithmetic and path meaning}
If the physical model declares constant-rate joint execution, $s$ can
parameterize its realized path after normalizing duration. If only endpoints
are declared, the same segment is an accounting interpolation. The endpoint
identity is exact in both uses. Calling it an observed physical history
requires the first additional declaration and, outside a simulation,
appropriate measurement.

More generally, a smooth potential integrates to the same total along any
admissible sufficiently regular path between the same endpoints. This does
not imply equal physical cost for every route: process burdens and other
state coordinates may differ. It also does not imply equal child allocations
when a group path is decomposed in different ways.

\section{A schedule is a function, not a promise about one quantity}
A pre-action quote can commit $q\mapsto E(q)$ on a physically accepted
domain. For the Ridge--Vale example the schedule is $4q-q^2/4$.
If the accepted upper quantity is four and the measured quantity is three,
the committed schedule gives $E(3)=39/4$, not twelve. The measurement selects
a point of an existing function; it does not authorize a new potential
chosen after the outcome.

Quantity shortfall does not always mean a smaller signed value. On a
decreasing part of a concave schedule, less movement can have greater value.
The schedule, not a verbal monotonicity assumption, governs the account.
If quantity exceeds the committed domain, the old positive quote cannot
simply be extended. The historical foundation retains an explicit
overexecution boundary rather than inventing punitive physics.

\section{Convexity and activation burden}
For a convex differentiable $V$,
$V(z+\delta)\geq V(z)+\nabla V(z)^T\delta$. Hence exact value is no
greater than its tangent value when both subtract the same $C$. For the
positive definite quadratic the inequality is strict for nonzero $\delta$.

Along a fixed edge, the field-benefit schedule is concave. The complete
schedule remains concave when the declared $C(q)$ is convex. Nonnegative
cost alone is not sufficient. A fixed activation cost with $C(0)=0$ and
$C(q)=c_0>0$ for $q>0$ creates a downward jump at zero; interpolation
across it would invent an unregistered discount. The first prospective
lossless domain avoids this complication by declaring $C=0$, not by
claiming real operations are universally costless.

\section{What a quote does not guarantee}
A valid schedule does not establish permission, ownership, sensor integrity,
or execution. A positive quote does not prove the action will occur. A
negative quote does not prove that essential service should be denied.
The mathematics describes represented field effect. Institutions and the
experimental capacity policy are additional objects with additional claims.

\section*{Chapter recap}
The exact finite difference and the integrated force are two expressions for
one field change. A committed schedule binds that expression to a baseline,
rule and domain. Physical-path meaning and institutional consequences do
not follow from integration alone.
